\documentclass[10pt,a4paper]{amsart}
\usepackage{amsmath,amssymb,amsthm,graphicx}
\usepackage{verbatim}
\newenvironment{reference}{\comment}{\endcomment}
%\newenvironment{reference}{}{}
\newenvironment{slogan}{\comment}{\endcomment}
\newenvironment{history}{\comment}{\endcomment}
% For commutative diagrams we use Xy-pic
\usepackage[all]{xy}

% We use 2cell for 2-commutative diagrams.
\xyoption{2cell}
% Theorem environments.
%
\theoremstyle{plain}
\newtheorem{theorem}[subsection]{Theorem}
\newtheorem{proposition}[subsection]{Proposition}
\newtheorem{lemma}[subsection]{Lemma}

\theoremstyle{definition}
\newtheorem{definition}[subsection]{Definition}
\newtheorem{example}[subsection]{Example}
\newtheorem{exercise}[subsection]{Exercise}
\newtheorem{situation}[subsection]{Situation}

\theoremstyle{remark}
\newtheorem{remark}[subsection]{Remark}
\newtheorem{remarks}[subsection]{Remarks}

\numberwithin{equation}{subsection}

% Macros
%
\def\lim{\mathop{\mathrm{lim}}\nolimits}
\def\colim{\mathop{\mathrm{colim}}\nolimits}
\def\Spec{\mathop{\mathrm{Spec}}}
\def\Hom{\mathop{\mathrm{Hom}}\nolimits}
\def\Ext{\mathop{\mathrm{Ext}}\nolimits}
\def\SheafHom{\mathop{\mathcal{H}\!\mathit{om}}\nolimits}
\def\SheafExt{\mathop{\mathcal{E}\!\mathit{xt}}\nolimits}
\def\Sch{\mathit{Sch}}
\def\Mor{\mathop{\mathrm{Mor}}\nolimits}
\def\Ob{\mathop{\mathrm{Ob}}\nolimits}
\def\Sh{\mathop{\mathit{Sh}}\nolimits}
\def\NL{\mathop{N\!L}\nolimits}
\def\CH{\mathop{\mathrm{CH}}\nolimits}
\def\proetale{{pro\text{-}\acute{e}tale}}
\def\etale{{\acute{e}tale}}
\def\QCoh{\mathit{QCoh}}
\def\Ker{\mathop{\mathrm{Ker}}}
\def\Im{\mathop{\mathrm{Im}}}
\def\Coker{\mathop{\mathrm{Coker}}}
\def\Coim{\mathop{\mathrm{Coim}}}

% Boxtimes
%
\DeclareMathSymbol{\boxtimes}{\mathbin}{AMSa}{"02}

%
% Macros for moduli stacks/spaces
%
\def\QCohstack{\mathcal{QC}\!\mathit{oh}}
\def\Cohstack{\mathcal{C}\!\mathit{oh}}
\def\Spacesstack{\mathcal{S}\!\mathit{paces}}
\def\Quotfunctor{\mathrm{Quot}}
\def\Hilbfunctor{\mathrm{Hilb}}
\def\Curvesstack{\mathcal{C}\!\mathit{urves}}
\def\Polarizedstack{\mathcal{P}\!\mathit{olarized}}
\def\Complexesstack{\mathcal{C}\!\mathit{omplexes}}
% \Pic is the operator that assigns to X its picard group, usage \Pic(X)
% \Picardstack_{X/B} denotes the Picard stack of X over B
% \Picardfunctor_{X/B} denotes the Picard functor of X over B
\def\Pic{\mathop{\mathrm{Pic}}\nolimits}
\def\Picardstack{\mathcal{P}\!\mathit{ic}}
\def\Picardfunctor{\mathrm{Pic}}
\def\Deformationcategory{\mathcal{D}\!\mathit{ef}}
\usepackage{cleveref}
\crefname{originaldefinition}{original definition}{original definitions}
\Crefname{originaldefinition}{Original definition}{Original definitions}
\makeatletter
\newlabel{definition-gysin-homomorphism}{{02T8}{source}{}{sourcepointer-1}{}}
\newlabel{definition-gysin-homomorphism@cref}{{[originaldefinition][0][]02T8}{[1][0][]source}{}{}{}}
\makeatother
\crefname{originallemma}{original lemma}{original lemmas}
\Crefname{originallemma}{Original lemma}{Original lemmas}
\makeatletter
\newlabel{lemma-equal-c1-as-cycles}{{02ST}{source}{}{sourcepointer-2}{}}
\newlabel{lemma-equal-c1-as-cycles@cref}{{[originallemma][0][]02ST}{[1][0][]source}{}{}{}}
\makeatother
\crefname{originalremark}{original remark}{original remarks}
\Crefname{originalremark}{Original remark}{Original remarks}
\makeatletter
\newlabel{remark-infinite-sums-rational-equivalences}{{02RZ}{source}{}{sourcepointer-3}{}}
\newlabel{remark-infinite-sums-rational-equivalences@cref}{{[originalremark][0][]02RZ}{[1][0][]source}{}{}{}}
\makeatother
\crefname{originalremark}{original remark}{original remarks}
\Crefname{originalremark}{Original remark}{Original remarks}
\makeatletter
\newlabel{remark-pullback-pairs}{{0B6Y}{source}{}{sourcepointer-4}{}}
\newlabel{remark-pullback-pairs@cref}{{[originalremark][0][]0B6Y}{[1][0][]source}{}{}{}}
\makeatother
\crefname{originalreference}{original reference}{original references}
\Crefname{originalreference}{Original reference}{Original references}
\makeatletter
\newlabel{situation-setup}{{02QL}{source}{}{sourcepointer-5}{}}
\newlabel{situation-setup@cref}{{[originalreference][0][]02QL}{[1][0][]source}{}{}{}}
\makeatother
\title{Compatibility of the Gysin map with a proper pushforward of the ambient scheme}
\author{The Stacks Project Authors}
\date{stacks-project:chow.tex:local-snapshot-20261012; GNU FDL 1.2 or later}
\begin{document}
\maketitle
\noindent\textit{Editorial note.} All delivered mathematical text is the author's own, copied from the pinned Stacks Project chapter file chow.tex; no statement or proof step was rewritten, reordered or removed.\par\medskip
\noindent\textit{Editorial note.} Preamble provenance: the pinned chapter file opens with \textbackslash{}input\{preamble\} and carries no preamble text of its own, so the project's single Stacks-wide preamble staged at sources/stacks-chow/preamble.tex is excerpted verbatim as three spans, lines 8-12 (the verbatim package together with the project's own reference, slogan and history display-environment definitions, including the author's commented-out older form of the reference environment), lines 14-18 (Xy-pic setup) and lines 154-217 (the native theorem declarations and mathematical macros). The lines 8-12 span is required because a counted statement or its surrounding text may use those display environments, which the project typesets as comments. Nothing was added, removed or redefined; the wrapper class and the amsmath/amssymb/amsthm/graphicx base packages are supplied by the assembler.\par\medskip
\noindent\textit{Editorial note.} Same-chapter pointers are located at their \textbackslash{}label lines inside the pinned chapter file; cross-chapter pointers are located in the project's own tag list staged in the snapshot (sources/stacks-chow/tags.txt), at the verbatim lines carrying their canonical tags.\par\medskip
\noindent\textit{Allowed standard prerequisites (editorial).}\begin{itemize}
\item The comparison of first Chern classes as cycles on integral schemes, and infinite sums of rational equivalences. Same chapter chow.tex, labels lemma-equal-c1-as-cycles (proof lines 5293 and 5297) and remark-infinite-sums-rational-equivalences (proof line 5298); the pullback diagram is described in remark-pullback-pairs (proof line 5277). The first identifies the two cycle classes in the two cases according to whether the component lies in the divisor, and the second extends the resulting equality from a single integral class to a general cycle.
\end{itemize}\medskip
\ifdefined\hypertarget\hypertarget{sourcepointer-1}{}\fi
\noindent\textit{Source reference.} Original definition 02T8: sources/stacks-chow/pinned-original-source.tex, source lines 5069--5069. Navigation pointer for \textbackslash{}ref\{definition-gysin-homomorphism\} used inside the counted span (line 5269); the argument that consumes it is retained in the span.\par\medskip
\ifdefined\hypertarget\hypertarget{sourcepointer-2}{}\fi
\noindent\textit{Source reference.} Original lemma 02ST: sources/stacks-chow/pinned-original-source.tex, source lines 4590--4590. Navigation pointer for \textbackslash{}ref\{lemma-equal-c1-as-cycles\} used inside the counted span (line 5293); the argument that consumes it is retained in the span.\par\medskip
\ifdefined\hypertarget\hypertarget{sourcepointer-3}{}\fi
\noindent\textit{Source reference.} Original remark 02RZ: sources/stacks-chow/pinned-original-source.tex, source lines 3057--3057. Navigation pointer for \textbackslash{}ref\{remark-infinite-sums-rational-equivalences\} used inside the counted span (line 5298); the argument that consumes it is retained in the span.\par\medskip
\ifdefined\hypertarget\hypertarget{sourcepointer-4}{}\fi
\noindent\textit{Source reference.} Original remark 0B6Y: sources/stacks-chow/pinned-original-source.tex, source lines 5248--5248. Navigation pointer for \textbackslash{}ref\{remark-pullback-pairs\} used inside the counted span (line 5277); the argument that consumes it is retained in the span.\par\medskip
\ifdefined\hypertarget\hypertarget{sourcepointer-5}{}\fi
\noindent\textit{Source reference.} Original reference 02QL: sources/stacks-chow/pinned-original-source.tex, source lines 1372--1372. Navigation pointer for \textbackslash{}ref\{situation-setup\} used inside the counted span (line 5265); the argument that consumes it is retained in the span.\par\medskip
\begin{lemma}
\label{lemma-closed-in-X-gysin}
Let $(S, \delta)$ be as in Situation \ref{situation-setup}.
Let $f : X' \to X$ be a proper morphism of schemes
locally of finite type over $S$.
Let $(\mathcal{L}, s, i : D \to X)$ be as in
Definition \ref{definition-gysin-homomorphism}.
Form the diagram
$$
\xymatrix{
D' \ar[d]_g \ar[r]_{i'} & X' \ar[d]^f \\
D \ar[r]^i & X
}
$$
as in Remark \ref{remark-pullback-pairs}.
For any $(k + 1)$-cycle $\alpha'$ on $X'$ we have
$i^*f_*\alpha' = g_*(i')^*\alpha'$ in $\CH_k(D)$
(this makes sense as $f_*$ is defined on the level of cycles).
\end{lemma}

\begin{proof}
Suppose $\alpha = [W']$ for some integral closed subscheme
$W' \subset X'$. Let $W = f(W') \subset X$. In case $W' \not \subset D'$,
then $W \not \subset D$ and we see that
$$
[W' \cap D']_k = \text{div}_{\mathcal{L}'|_{W'}}({s'|_{W'}})
\quad\text{and}\quad
[W \cap D]_k = \text{div}_{\mathcal{L}|_W}(s|_W)
$$
and hence $f_*$ of the first cycle equals the second cycle by
Lemma \ref{lemma-equal-c1-as-cycles}. Hence the
equality holds as cycles. In case $W' \subset D'$, then
$W \subset D$ and $f_*(c_1(\mathcal{L}|_{W'}) \cap [W'])$
is equal to $c_1(\mathcal{L}|_W) \cap [W]$ in $\CH_k(W)$ by the second
assertion of Lemma \ref{lemma-equal-c1-as-cycles}.
By Remark \ref{remark-infinite-sums-rational-equivalences}
the result follows for general $\alpha'$.
\end{proof}

\end{document}
